%%%%%%%%%%%%%%%%%%%%%%%%%%%%%%%%%%%%%%%%%%%%%%%%%%%%%%%%%%%%%%%%%%%
% Physik IV - Aufgabenblatt 2: Ausarbeitung, Tafelvortrag, Lernfassung
% Standalone LaTeX file
%%%%%%%%%%%%%%%%%%%%%%%%%%%%%%%%%%%%%%%%%%%%%%%%%%%%%%%%%%%%%%%%%%%

\documentclass[11pt,a4paper]{scrartcl}

%%%%%%%%%%%%%%%%%%%%%%%%%%%%%%%%%%%%%%%%%%%%%%%%%%%%%%%%%%%%%%%%%%%
% Metadata
%%%%%%%%%%%%%%%%%%%%%%%%%%%%%%%%%%%%%%%%%%%%%%%%%%%%%%%%%%%%%%%%%%%
\newcommand{\CourseTitle}{Proseminar zur Physik IV: Kerne und Teilchen}
\newcommand{\CourseShortTitle}{Physik IV}
\newcommand{\DocumentKind}{Hausaufgabe / Tafelvortrag / Lernfassung}
\newcommand{\DocumentTitle}{Aufgabenblatt 2}
\newcommand{\DocumentSubtitle}{Formfaktoren und Grenzen der Rutherford-Streuung}
\newcommand{\AuthorName}{Tobias Laser}
\newcommand{\InstitutionName}{Universit\"at Innsbruck}
\newcommand{\SubmissionDate}{\today}

\newcommand{\makeuniversitytitle}{%
    \begin{titlepage}
        \centering
        \vspace*{2cm}
        {\Large \CourseTitle\par}
        \vspace{1.2cm}
        {\huge\bfseries \DocumentTitle\par}
        \vspace{0.4cm}
        {\Large \DocumentSubtitle\par}
        \vspace{1cm}
        {\Large \DocumentKind\par}
        \vfill
        {\large \AuthorName\par}
        \vspace{0.2cm}
        {\large \InstitutionName\par}
        \vspace{0.8cm}
        {\large \SubmissionDate\par}
    \end{titlepage}
    \pagenumbering{arabic}
}

%%%%%%%%%%%%%%%%%%%%%%%%%%%%%%%%%%%%%%%%%%%%%%%%%%%%%%%%%%%%%%%%%%%
% Packages and setup
%%%%%%%%%%%%%%%%%%%%%%%%%%%%%%%%%%%%%%%%%%%%%%%%%%%%%%%%%%%%%%%%%%%
\usepackage[margin=2.5cm]{geometry}
\usepackage[onehalfspacing]{setspace}
\usepackage{scrlayer-scrpage}
\usepackage[T1]{fontenc}
\usepackage[utf8]{inputenc}
\usepackage[ngerman]{babel}
\usepackage{lmodern}
\usepackage{microtype}
\usepackage{csquotes}
\usepackage{amsmath,amsthm,amssymb,mathtools}
\usepackage{bm}
\usepackage{icomma}
\usepackage{cancel}
\usepackage{graphicx}
\usepackage{tikz}
\usetikzlibrary{arrows.meta,calc,decorations.pathmorphing,positioning,patterns,angles,quotes}
\usepackage{pgfplots}
\pgfplotsset{compat=1.18}
\usepackage[locale=DE,space-before-unit=true,per-mode=symbol,separate-uncertainty=true]{siunitx}
\usepackage{booktabs,multirow,array,tabularx,longtable}
\usepackage{caption,subcaption}
\usepackage{placeins}
\usepackage{enumitem}
\usepackage{xparse}
\usepackage[most]{tcolorbox}
\tcbuselibrary{breakable,skins,theorems}
\usepackage[breaklinks=true,colorlinks=true,linkcolor=blue,urlcolor=blue,citecolor=blue]{hyperref}
\usepackage[nameinlink,noabbrev]{cleveref}

\clearpairofpagestyles
\ihead{\CourseShortTitle}
\ohead{\DocumentKind}
\cfoot{\pagemark}
\pagestyle{scrheadings}

\setlist[enumerate]{itemsep=0.4em,topsep=0.4em}
\setlist[itemize]{itemsep=0.25em,topsep=0.35em}
\captionsetup{font=small,labelfont=bf}

%%%%%%%%%%%%%%%%%%%%%%%%%%%%%%%%%%%%%%%%%%%%%%%%%%%%%%%%%%%%%%%%%%%
% Macros
%%%%%%%%%%%%%%%%%%%%%%%%%%%%%%%%%%%%%%%%%%%%%%%%%%%%%%%%%%%%%%%%%%%
\newcommand{\R}{\mathbb{R}}
\newcommand{\abs}[1]{\left| #1 \right|}
\newcommand{\norm}[1]{\left\lVert #1 \right\rVert}
\newcommand{\paren}[1]{\left( #1 \right)}
\newcommand{\bracket}[1]{\left[ #1 \right]}
\newcommand{\dd}{\,\mathrm{d}}
\newcommand{\dv}[2]{\frac{\dd #1}{\dd #2}}
\newcommand{\vect}[1]{\bm{#1}}
\newcommand{\unitvect}[1]{\hat{\bm{#1}}}
\newcommand{\half}{\frac{1}{2}}
\newcommand{\hbarc}{\hbar c}
\newcommand{\MeV}{\si{MeV}}
\newcommand{\fm}{\si{fm}}
\newcommand{\nuclide}[3]{^{#1}_{#2}\mathrm{#3}}
\DeclareSIUnit{\fermi}{fm}

%%%%%%%%%%%%%%%%%%%%%%%%%%%%%%%%%%%%%%%%%%%%%%%%%%%%%%%%%%%%%%%%%%%
% Boxes
%%%%%%%%%%%%%%%%%%%%%%%%%%%%%%%%%%%%%%%%%%%%%%%%%%%%%%%%%%%%%%%%%%%
\tcbset{
    unibox/.style={
        enhanced,
        breakable,
        sharp corners,
        boxrule=0.6pt,
        left=1.2mm,
        right=1.2mm,
        top=1mm,
        bottom=1mm,
        before skip=0.8em,
        after skip=0.8em,
        fonttitle=\bfseries
    }
}

\newtcolorbox{taskbox}[1][]{unibox,title=Aufgabe,colback=white,colframe=black!75,#1}
\newtcolorbox{calculationbox}[1][]{unibox,title=Rechnung,colback=black!2,colframe=black!55,#1}
\newtcolorbox{sidecalcbox}[1][]{unibox,title=Nebenrechnung,colback=black!1,colframe=black!40,#1}
\newtcolorbox{solutionbox}[1][]{unibox,title=L\"osung,colback=blue!2,colframe=blue!45!black,#1}
\newtcolorbox{answerbox}[1][]{unibox,title=Antwort,colback=green!3,colframe=green!45!black,#1}
\newtcolorbox{notebox}[1][]{unibox,title=Notiz,colback=yellow!6,colframe=yellow!45!black,#1}
\newtcolorbox{definitionbox}[1][]{unibox,title=Definition,colback=cyan!3,colframe=cyan!45!black,#1}
\newtcolorbox{warningbox}[1][]{unibox,title=Achtung,colback=red!3,colframe=red!55!black,#1}
\newtcolorbox{summarybox}[1][]{unibox,title=Zusammenfassung,colback=violet!3,colframe=violet!50!black,#1}
\newtcolorbox{formulabox}[1][]{unibox,title=Formel,colback=orange!4,colframe=orange!65!black,#1}

%%%%%%%%%%%%%%%%%%%%%%%%%%%%%%%%%%%%%%%%%%%%%%%%%%%%%%%%%%%%%%%%%%%
\begin{document}

\makeuniversitytitle
\tableofcontents
\newpage

%%%%%%%%%%%%%%%%%%%%%%%%%%%%%%%%%%%%%%%%%%%%%%%%%%%%%%%%%%%%%%%%%%%
\section*{Arbeitsnotiz}
\addcontentsline{toc}{section}{Arbeitsnotiz}

Diese Ausarbeitung ist bewusst nicht nur eine knappe Abgabe, sondern auch eine Lernfassung. Die Musterl\"osungen wurden nur dort verwendet, wo die Schritte nachvollziehbar und sinnvoll waren; die Herleitungen sind hier sauber und geschlossen neu formuliert. Die beiden Kernthemen sind:
\begin{itemize}
    \item der Formfaktor als Fourier-Transformierte einer Ladungsverteilung,
    \item die Grenze der Rutherford-Streuung durch endliche Kerngr\"o\ss{}e und Kernkr\"afte.
\end{itemize}

\begin{summarybox}[title=Endergebnisse auf einen Blick]
\begin{align*}
    F(q) &= 4\pi\int_0^\infty \frac{\sin(qr/\hbar)}{qr/\hbar} f(r)r^2\dd r,\\
    F_{\text{hom. Kugel}}(q) &= 3\frac{\sin u-u\cos u}{u^3}, \qquad u=\frac{qR}{\hbar},\\
    F(0)&=1,\\[0.4em]
    \delta_0(E=\SI{10}{MeV})&=\SI{47.2}{fm},\qquad
    \delta_0(E=\SI{80}{MeV})=\SI{5.90}{fm},\\
    R_{\mathrm{Pb}}&=\SI{7.70}{fm},\\
    \delta(\theta=\SI{40}{\degree},\SI{10}{MeV})&=\SI{92.7}{fm},\\
    \delta(\theta=\SI{40}{\degree},\SI{80}{MeV})&=\SI{11.6}{fm},\\
    \theta_{\mathrm{krit}}(\SI{80}{MeV})&\approx \SI{76.8}{\degree}.
\end{align*}
\end{summarybox}

\newpage

%%%%%%%%%%%%%%%%%%%%%%%%%%%%%%%%%%%%%%%%%%%%%%%%%%%%%%%%%%%%%%%%%%%
\section{Aufgabe 1: Formfaktoren -- Herleitung}

\begin{taskbox}
In der Vorlesung wurde der Formfaktor einer Ladungsverteilung $f(\vect{x})$ definiert als
\[
    F(\vect{q})=\int e^{i\vect{q}\cdot\vect{x}/\hbar}f(\vect{x})\dd^3x.
\]
F\"ur kugelsymmetrische Ladungsverteilungen soll gezeigt werden, dass $F$ nur vom Betrag $q=\abs{\vect{q}}$ abh\"angt und sich auf ein eindimensionales Integral reduzieren l\"asst. Danach soll der Spezialfall einer homogen geladenen Kugel berechnet werden.
\end{taskbox}

\subsection{Teil (a): Reduktion auf ein radiales Integral}

\begin{definitionbox}[title=Physikalische Idee]
Der Formfaktor beschreibt, wie stark eine ausgedehnte Ladungsverteilung bei einem Impuls\"ubertrag $\vect{q}$ streut. Mathematisch ist er im Wesentlichen die Fourier-Transformierte der Ladungsverteilung. Bei Kugelsymmetrie kann es keine bevorzugte Raumrichtung geben; deshalb darf $F$ nur von $q=\abs{\vect{q}}$ abh\"angen.
\end{definitionbox}

\begin{calculationbox}[title=Koordinatenwahl und Skalarprodukt]
Da $f(\vect{x})=f(r)$ kugelsymmetrisch ist, darf man die $z$-Achse in Richtung von $\vect{q}$ legen:
\[
    \vect{q}=q\unitvect{e}_z=\begin{pmatrix}0\\0\\q\end{pmatrix}.
\]
In Kugelkoordinaten gilt
\[
    \vect{x}=\begin{pmatrix}
        r\sin\theta\cos\varphi\\
        r\sin\theta\sin\varphi\\
        r\cos\theta
    \end{pmatrix},
    \qquad
    \vect{q}\cdot\vect{x}=qr\cos\theta.
\]
Das Volumenelement ist
\[
    \dd^3x=r^2\sin\theta\dd r\dd\theta\dd\varphi.
\]
Damit wird
\[
    F(q)=\int_0^\infty\int_0^\pi\int_0^{2\pi}
    e^{iqr\cos\theta/\hbar}f(r)r^2\sin\theta
    \dd\varphi\dd\theta\dd r.
\]
\end{calculationbox}

\begin{calculationbox}[title=Winkelintegrale]
Das Integral \"uber $\varphi$ gibt direkt den Faktor $2\pi$:
\[
    F(q)=2\pi\int_0^\infty f(r)r^2
    \bracket{\int_0^\pi e^{iqr\cos\theta/\hbar}\sin\theta\dd\theta}\dd r.
\]
F\"ur das $\theta$-Integral setzen wir
\[
    u=\cos\theta,\qquad \dd u=-\sin\theta\dd\theta.
\]
Die Grenzen werden $\theta=0\mapsto u=1$ und $\theta=\pi\mapsto u=-1$. Also
\begin{align*}
    \int_0^\pi e^{iqr\cos\theta/\hbar}\sin\theta\dd\theta
    &=\int_{1}^{-1}e^{iqru/\hbar}(-\dd u)\\
    &=\int_{-1}^{1}e^{iqru/\hbar}\dd u\\
    &=\left.\frac{e^{iqru/\hbar}}{iqr/\hbar}\right|_{-1}^{1}\\
    &=\frac{e^{iqr/\hbar}-e^{-iqr/\hbar}}{iqr/\hbar}\\
    &=2\frac{\sin(qr/\hbar)}{qr/\hbar}.
\end{align*}
\end{calculationbox}

\begin{answerbox}
Damit folgt
\[
    \boxed{F(q)=4\pi\int_0^\infty \frac{\sin(qr/\hbar)}{qr/\hbar}f(r)r^2\dd r.}
\]
Der Formfaktor h\"angt nur von $q=\abs{\vect{q}}$ ab, weil alle Richtungsinformationen durch die Kugelsymmetrie herausintegriert werden.
\end{answerbox}

\subsection{Teil (b): Homogen geladene Kugel}

\begin{definitionbox}[title=Ladungsverteilung einer homogen geladenen Kugel]
Eine homogen geladene Kugel mit Radius $R$ hat konstante Ladungsdichte im Inneren und keine Ladung au\ss{}erhalb:
\[
    f(r)=\begin{cases}
        C, & 0\le r\le R,\\
        0, & r>R.
    \end{cases}
\]
Die Aufgabenstellung verwendet die Normierung
\[
    \int f(\vect{x})\dd^3x=1.
\]
\end{definitionbox}

\begin{calculationbox}[title=Normierung der Ladungsdichte]
Aus der Normierung folgt
\begin{align*}
    1&=\int f(\vect{x})\dd^3x
      =4\pi\int_0^\infty f(r)r^2\dd r
      =4\pi\int_0^R Cr^2\dd r\\
     &=4\pi C\left.\frac{r^3}{3}\right|_0^R
      =\frac{4\pi CR^3}{3}.
\end{align*}
Damit ist
\[
    C=\frac{3}{4\pi R^3}.
\]
Also
\[
    f(r)=\begin{cases}
        \dfrac{3}{4\pi R^3}, & 0\le r\le R,\\[0.4em]
        0, & r>R.
    \end{cases}
\]
\end{calculationbox}

\begin{calculationbox}[title=Einsetzen in den Formfaktor]
Wir setzen die Dichte in das Ergebnis aus Teil (a) ein:
\begin{align*}
    F(q)&=4\pi\int_0^R \frac{\sin(qr/\hbar)}{qr/\hbar}\frac{3}{4\pi R^3}r^2\dd r\\
        &=\frac{3}{R^3}\int_0^R \frac{\sin(qr/\hbar)}{qr/\hbar}r^2\dd r.
\end{align*}
Mit $a=q/\hbar$ wird
\[
    F(q)=\frac{3}{R^3a}\int_0^R r\sin(ar)\dd r.
\]
Das verbleibende Integral ist
\begin{align*}
    \int r\sin(ar)\dd r
    &=-\frac{r\cos(ar)}{a}+\frac{\sin(ar)}{a^2}.
\end{align*}
Somit
\begin{align*}
    F(q)&=\frac{3}{R^3a}\left[-\frac{r\cos(ar)}{a}+\frac{\sin(ar)}{a^2}\right]_0^R\\
        &=\frac{3}{R^3a}\left[-\frac{R\cos(aR)}{a}+\frac{\sin(aR)}{a^2}\right]\\
        &=3\frac{\sin(aR)-aR\cos(aR)}{(aR)^3}.
\end{align*}
\end{calculationbox}

\begin{answerbox}
Mit
\[
    \boxed{u=\frac{qR}{\hbar}}
\]
erh\"alt man
\[
    \boxed{F(q)=3\frac{\sin u-u\cos u}{u^3}.}
\]
Eine homogen geladene Kugel ist eine brauchbare erste N\"aherung f\"ur mittlere und schwere Atomkerne. Realistischer ist f\"ur Kerne jedoch eine ann\"ahernd konstante Innendichte mit diffuser Randzone.
\end{answerbox}

\subsection{Teil (c): Formfaktor bei $q=0$}

\begin{solutionbox}
Man kann entweder direkt in die urspr\"ungliche Definition einsetzen oder den Grenzwert des radialen Integrals nehmen. Direkt aus der Definition folgt
\[
    F(0)=\int e^{i\cdot 0}f(\vect{x})\dd^3x
    =\int f(\vect{x})\dd^3x.
\]
Da $f$ normiert ist, gilt
\[
    \int f(\vect{x})\dd^3x=1.
\]
\end{solutionbox}

\begin{answerbox}
\[
    \boxed{F(0)=1.}
\]
Physikalisch bedeutet das: Bei verschwindendem Impuls\"ubertrag wird die Struktur nicht aufgel\"ost; man sieht nur die gesamte normierte Ladung.
\end{answerbox}

\begin{warningbox}[title=Typische Fehler bei Aufgabe 1]
\begin{itemize}
    \item Nicht $q=0$ direkt in $\sin(qr/\hbar)/(qr/\hbar)$ als $0/0$ behandeln. Der Grenzwert ist $1$.
    \item Beim Winkelintegral das Minuszeichen der Substitution $u=\cos\theta$ korrekt mit den Integrationsgrenzen verrechnen.
    \item Bei der homogenen Kugel nicht vergessen, die Dichte zuerst zu normieren.
    \item $F(q)$ ist dimensionslos; deshalb muss $u=qR/\hbar$ dimensionslos sein.
\end{itemize}
\end{warningbox}

\FloatBarrier
\newpage

%%%%%%%%%%%%%%%%%%%%%%%%%%%%%%%%%%%%%%%%%%%%%%%%%%%%%%%%%%%%%%%%%%%
\section{Aufgabe 2: Grenzen der Rutherford-Streuung}

\begin{taskbox}
Der Rutherford-Streuquerschnitt beschreibt eine Streuung korrekt, solange die Teilchen einander nicht so nahe kommen, dass Kernkr\"afte signifikant werden. Gegeben ist der Zusammenhang zwischen Streuwinkel $\theta$ und kleinstem Abstand $\delta$:
\[
    \delta=\frac{\delta_0}{2}\paren{1+\frac{1}{\sin(\theta/2)}}.
\]
Untersucht wird die Streuung eines $\nuclide{7}{4}{Be}$-Kerns an einem $\nuclide{208}{82}{Pb}$-Kern f\"ur $E_{\mathrm{kin}}=\SI{10}{MeV}$ und $\SI{80}{MeV}$.
\end{taskbox}

\subsection{Vorbereitung: Coulomb-Konstante in Kernphysik-Einheiten}

\begin{formulabox}
F\"ur zwei Kerne mit Kernladungszahlen $Z_1$ und $Z_2$ gilt f\"ur die Coulomb-Energie im Abstand $r$:
\[
    E_{\mathrm{pot}}(r)=\frac{1}{4\pi\varepsilon_0}\frac{Z_1Z_2e^2}{r}.
\]
In Kernphysik-Einheiten verwendet man praktisch
\[
    \frac{e^2}{4\pi\varepsilon_0}=\SI{1.44}{MeV.fm}.
\]
\end{formulabox}

\subsection{Teil (a): Bestimmung von $\delta_0$}

\begin{calculationbox}[title=Zentraler Sto\ss{}]
Beim zentralen Sto\ss{} ist
\[
    \theta=\SI{180}{\degree}=\pi.
\]
Damit ist
\[
    \sin\paren{\frac{\theta}{2}}=\sin\paren{\frac{\pi}{2}}=1.
\]
Einsetzen in die gegebene Abstandsgleichung liefert
\[
    \delta=\frac{\delta_0}{2}(1+1)=\delta_0.
\]
Also ist $\delta_0$ genau der kleinste Abstand beim zentralen Sto\ss{}.
\end{calculationbox}

\begin{calculationbox}[title=Energieumwandlung bei gr\"o\ss{}ter Ann\"aherung]
Beim zentralen Sto\ss{} wird am Umkehrpunkt die gesamte kinetische Energie in Coulomb-Potentialenergie umgewandelt:
\[
    E_{\mathrm{kin}}=\frac{1}{4\pi\varepsilon_0}\frac{Z_1Z_2e^2}{\delta_0}.
\]
Daraus folgt
\[
    \delta_0=\frac{1}{4\pi\varepsilon_0}\frac{Z_1Z_2e^2}{E_{\mathrm{kin}}}.
\]
F\"ur $\nuclide{7}{4}{Be}$ an $\nuclide{208}{82}{Pb}$ gilt
\[
    Z_1=4,\qquad Z_2=82,
\]
also
\[
    Z_1Z_2=328.
\]
Damit
\[
    \delta_0=\frac{328\cdot\SI{1.44}{MeV.fm}}{E_{\mathrm{kin}}}
    =\frac{\SI{472.32}{MeV.fm}}{E_{\mathrm{kin}}}.
\]
\end{calculationbox}

\begin{answerbox}
\[
    \boxed{\delta_0(\SI{10}{MeV})=\frac{\SI{472.32}{MeV.fm}}{\SI{10}{MeV}}=\SI{47.2}{fm}}
\]
\[
    \boxed{\delta_0(\SI{80}{MeV})=\frac{\SI{472.32}{MeV.fm}}{\SI{80}{MeV}}=\SI{5.90}{fm}}
\]
Je h\"oher die kinetische Energie ist, desto kleiner wird der kleinste Abstand beim zentralen Sto\ss{}.
\end{answerbox}

\subsection{Teil (b): Vergleich mit dem Radius eines Bleikerns}

\begin{calculationbox}
Mit der Kernradius-Formel
\[
    R=r_0A^{1/3}
\]
und
\[
    r_0=\SI{1.3}{fm},\qquad A=208
\]
erh\"alt man
\[
    R_{\mathrm{Pb}}=\SI{1.3}{fm}\cdot 208^{1/3}=\SI{7.70}{fm}.
\]
Der Durchmesser ist
\[
    D_{\mathrm{Pb}}=2R_{\mathrm{Pb}}=\SI{15.4}{fm}.
\]
\end{calculationbox}

\begin{answerbox}
\begin{center}
\begin{tabular}{@{}lcc@{}}
\toprule
Fall & Abstand $\delta_0$ & Vergleich mit $R_{\mathrm{Pb}}=\SI{7.70}{fm}$ \\
\midrule
$E_{\mathrm{kin}}=\SI{10}{MeV}$ & $\SI{47.2}{fm}$ & deutlich au\ss{}erhalb des Pb-Kerns \\
$E_{\mathrm{kin}}=\SI{80}{MeV}$ & $\SI{5.90}{fm}$ & kleiner als $R_{\mathrm{Pb}}$ \\
\bottomrule
\end{tabular}
\end{center}
Bei $\SI{10}{MeV}$ kommt das Projektil selbst beim zentralen Sto\ss{} nicht nahe genug an den Kern, um den Kerndurchmesser gut zu untersuchen. Bei $\SI{80}{MeV}$ kann das Projektil bei gro\ss{}en Ablenkwinkeln in den Bereich des Bleikerns kommen; dort sind Abweichungen von der reinen Rutherford-Streuung zu erwarten. Genau solche Abweichungen liefern Informationen \"uber Kerngr\"o\ss{}e und Kerndurchmesser.
\end{answerbox}

\subsection{Teil (c): Abstand bei $\theta=\SI{40}{\degree}$}

\begin{calculationbox}
Die Abstandsgleichung lautet
\[
    \delta(\theta)=\frac{\delta_0}{2}\paren{1+\frac{1}{\sin(\theta/2)}}.
\]
F\"ur $\theta=\SI{40}{\degree}$ ist
\[
    \sin\paren{\frac{\theta}{2}}=\sin(\SI{20}{\degree})\approx 0{,}342.
\]
Damit
\[
    \frac{1}{2}\paren{1+\frac{1}{\sin(\SI{20}{\degree})}}
    \approx 1{,}962.
\]
Also gilt
\[
    \delta(\SI{40}{\degree})\approx 1{,}962\,\delta_0.
\]
\end{calculationbox}

\begin{answerbox}
\[
    \boxed{\delta(\SI{40}{\degree},\SI{10}{MeV})
    \approx 1{,}962\cdot\SI{47.2}{fm}=\SI{92.7}{fm}}
\]
\[
    \boxed{\delta(\SI{40}{\degree},\SI{80}{MeV})
    \approx 1{,}962\cdot\SI{5.90}{fm}=\SI{11.6}{fm}}
\]
Beide Abst\"ande liegen bei $\theta=\SI{40}{\degree}$ noch oberhalb von $R_{\mathrm{Pb}}=\SI{7.70}{fm}$. Bei diesem Winkel ist daher noch keine starke Abweichung von Rutherford zu erwarten.
\end{answerbox}

\subsection{Teil (d): Kritischer Ablenkwinkel bei $E_{\mathrm{kin}}=\SI{80}{MeV}$}

\begin{calculationbox}[title=Bedingung f\"ur Abweichung]
Die Rutherford-Beschreibung beginnt zu versagen, wenn der kleinste Abstand in die Gr\"o\ss{}enordnung des Bleikernradius kommt. Gem\"a\ss{} Hinweis setzen wir
\[
    \delta(\theta_{\mathrm{krit}})=R_{\mathrm{Pb}}.
\]
F\"ur $E_{\mathrm{kin}}=\SI{80}{MeV}$ ist
\[
    \delta_0=\SI{5.904}{fm}.
\]
Damit
\begin{align*}
    R_{\mathrm{Pb}}&=\frac{\delta_0}{2}\paren{1+\frac{1}{\sin(\theta/2)}}\\
    \frac{2R_{\mathrm{Pb}}}{\delta_0}&=1+\frac{1}{\sin(\theta/2)}\\
    \frac{1}{\sin(\theta/2)}&=\frac{2R_{\mathrm{Pb}}}{\delta_0}-1\\
    \sin(\theta/2)&=\frac{1}{\frac{2R_{\mathrm{Pb}}}{\delta_0}-1}.
\end{align*}
Einsetzen ergibt
\[
    \sin(\theta/2)=\frac{1}{\frac{2\cdot\SI{7.70}{fm}}{\SI{5.904}{fm}}-1}
    \approx 0{,}621.
\]
Also
\[
    \theta=2\arcsin(0{,}621)\approx \SI{76.8}{\degree}.
\]
\end{calculationbox}

\begin{answerbox}
\[
    \boxed{\theta_{\mathrm{krit}}\approx \SI{76.8}{\degree}.}
\]
F\"ur $E_{\mathrm{kin}}=\SI{80}{MeV}$ sind ab Ablenkwinkeln von etwa $\SI{77}{\degree}$ Abweichungen vom Rutherford-Wirkungsquerschnitt zu erwarten. F\"ur gr\"o\ss{}ere Winkel wird der kleinste Abstand noch kleiner und die Kernkr\"afte k\"onnen st\"arker beitragen.
\end{answerbox}

\begin{calculationbox}[title=Skizze: kleinster Abstand als Funktion des Streuwinkels]
\begin{center}
\begin{tikzpicture}
\begin{axis}[
    width=0.92\textwidth,
    height=7cm,
    xlabel={Streuwinkel $\theta$ in Grad},
    ylabel={kleinster Abstand $\delta$ in fm},
    xmin=10, xmax=180,
    ymin=0, ymax=120,
    grid=both,
    legend pos=north east,
    samples=220
]
\addplot[thick,domain=10:180]{47.232/2*(1+1/sin(x/2))};
\addlegendentry{$E=\SI{10}{MeV}$}
\addplot[thick,dashed,domain=10:180]{5.904/2*(1+1/sin(x/2))};
\addlegendentry{$E=\SI{80}{MeV}$}
\addplot[thick,dotted,domain=10:180]{7.70249};
\addlegendentry{$R_{\mathrm{Pb}}$}
\end{axis}
\end{tikzpicture}
\end{center}
\end{calculationbox}

\begin{warningbox}[title=Typische Fehler bei Aufgabe 2]
\begin{itemize}
    \item Beim zentralen Sto\ss{} gilt $\theta=\SI{180}{\degree}$, nicht $\SI{0}{\degree}$.
    \item In der Gleichung steht $\sin(\theta/2)$, nicht $\sin\theta/2$ im Sinne von $(\sin\theta)/2$.
    \item Die Coulomb-Konstante in Kernphysik-Einheiten ist $e^2/(4\pi\varepsilon_0)=\SI{1.44}{MeV.fm}$.
    \item Der Radius $R=r_0A^{1/3}$ ist ein Radius, nicht der Durchmesser.
    \item Bei kleinen Winkeln ist $\delta$ gro\ss{}; gro\ss{}e Ablenkwinkel entsprechen kleinen Abst\"anden.
\end{itemize}
\end{warningbox}

\FloatBarrier
\newpage

%%%%%%%%%%%%%%%%%%%%%%%%%%%%%%%%%%%%%%%%%%%%%%%%%%%%%%%%%%%%%%%%%%%
\section{Tafelvortrag}

\subsection{Kurzstruktur f\"ur einen Vortrag von etwa 10--15 Minuten}

\begin{summarybox}[title=Empfohlene Tafelgliederung]
\begin{enumerate}[label=\arabic*.]
    \item \textbf{Motivation:} Formfaktor misst Struktur; Rutherford-Streuung gilt nur bei reinem Coulomb-Potential.
    \item \textbf{Aufgabe 1(a):} Achse entlang $\vect{q}$ w\"ahlen, Winkelintegral rechnen, radiales Integral zeigen.
    \item \textbf{Aufgabe 1(b):} homogene Kugel normieren, einsetzen, $u=qR/\hbar$ definieren.
    \item \textbf{Aufgabe 1(c):} $F(0)=1$ aus Normierung.
    \item \textbf{Aufgabe 2(a):} zentraler Sto\ss{}, $\delta_0$ aus Energieerhaltung.
    \item \textbf{Aufgabe 2(b)--(d):} Vergleich mit $R_{\mathrm{Pb}}$, kritischer Winkel.
    \item \textbf{Schlusssatz:} Abweichungen von Rutherford sind kein Fehler, sondern ein Messsignal f\"ur Kerngr\"o\ss{}e und Kernkr\"afte.
\end{enumerate}
\end{summarybox}

\subsection{Tafelbild 1: Formfaktor}

\begin{calculationbox}[title=Tafelbild]
Links oben:
\[
    F(\vect{q})=\int e^{i\vect{q}\cdot\vect{x}/\hbar}f(\vect{x})\dd^3x,
    \qquad f(\vect{x})=f(r).
\]
Dann direkt die Koordinatenwahl:
\[
    \vect{q}=q\unitvect{e}_z,
    \qquad \vect{q}\cdot\vect{x}=qr\cos\theta,
    \qquad \dd^3x=r^2\sin\theta\dd r\dd\theta\dd\varphi.
\]
Dann in einer Zeile das Ergebnis des Winkelintegrals:
\[
    \int_0^\pi e^{iqr\cos\theta/\hbar}\sin\theta\dd\theta
    =2\frac{\sin(qr/\hbar)}{qr/\hbar}.
\]
Gro\ss{} eingerahmt:
\[
    F(q)=4\pi\int_0^\infty\frac{\sin(qr/\hbar)}{qr/\hbar}f(r)r^2\dd r.
\]
\end{calculationbox}

\begin{notebox}[title=Sprechtext dazu]
\enquote{Ich darf die $z$-Achse entlang $\vect{q}$ legen, weil die Ladungsverteilung kugelsymmetrisch ist. Dadurch verschwindet die Richtungsabh\"angigkeit und nur der Betrag $q$ bleibt \"ubrig. Das ist genau der Grund, warum aus $F(\vect{q})$ ein $F(q)$ wird.}
\end{notebox}

\subsection{Tafelbild 2: Homogene Kugel}

\begin{calculationbox}[title=Tafelbild]
Zuerst die Dichte:
\[
    f(r)=\begin{cases}
        C, & r\le R,\\
        0, & r>R,
    \end{cases}
    \qquad
    1=4\pi C\frac{R^3}{3}
    \Rightarrow
    C=\frac{3}{4\pi R^3}.
\]
Dann einsetzen:
\[
    F(q)=\frac{3}{R^3}\int_0^R\frac{\sin(qr/\hbar)}{qr/\hbar}r^2\dd r.
\]
Mit $u=qR/\hbar$:
\[
    F(q)=3\frac{\sin u-u\cos u}{u^3}.
\]
\end{calculationbox}

\subsection{Tafelbild 3: Rutherford-Grenze}

\begin{calculationbox}[title=Tafelbild]
Grundgleichungen:
\[
    \delta=\frac{\delta_0}{2}\paren{1+\frac{1}{\sin(\theta/2)}}.
\]
Zentraler Sto\ss{}:
\[
    \theta=180^\circ\Rightarrow \delta=\delta_0.
\]
Energieerhaltung:
\[
    E_{\mathrm{kin}}=\frac{Z_1Z_2e^2}{4\pi\varepsilon_0\delta_0}
    \Rightarrow
    \delta_0=\frac{Z_1Z_2\cdot\SI{1.44}{MeV.fm}}{E_{\mathrm{kin}}}.
\]
Zahlen:
\[
    Z_1Z_2=4\cdot82=328,
    \qquad
    \delta_0=\frac{\SI{472.32}{MeV.fm}}{E_{\mathrm{kin}}}.
\]
\[
    \delta_0(10)=\SI{47.2}{fm},
    \qquad
    \delta_0(80)=\SI{5.90}{fm}.
\]
Bleiradius:
\[
    R_{\mathrm{Pb}}=\SI{1.3}{fm}\cdot 208^{1/3}=\SI{7.70}{fm}.
\]
Kritischer Winkel:
\[
    R_{\mathrm{Pb}}=\frac{\delta_0}{2}\paren{1+\frac{1}{\sin(\theta/2)}}
    \Rightarrow
    \theta\approx 76{,}8^\circ.
\]
\end{calculationbox}

\begin{notebox}[title=Sprechtext dazu]
\enquote{Bei $\SI{10}{MeV}$ bleibt das Beryllium selbst beim zentralen Sto\ss{} weit au\ss{}erhalb des Bleikerns. Bei $\SI{80}{MeV}$ ist der zentrale Abstand kleiner als der Bleiradius. Daher ist diese Energie geeignet, um \"uber Abweichungen von Rutherford Informationen \"uber die Kerngr\"o\ss{}e zu bekommen.}
\end{notebox}

\newpage

%%%%%%%%%%%%%%%%%%%%%%%%%%%%%%%%%%%%%%%%%%%%%%%%%%%%%%%%%%%%%%%%%%%
\section{Lernzusammenfassung}

\begin{summarybox}[title=Was man wirklich verstehen sollte]
\begin{itemize}
    \item Der Formfaktor ist die Fourier-Transformierte der Ladungsverteilung. Kleine $q$ bedeuten grobe Aufl\"osung, gro\ss{}e $q$ bedeuten feine Strukturaufl\"osung.
    \item Kugelsymmetrie reduziert das dreidimensionale Integral auf ein radiales Integral.
    \item $F(0)=1$ folgt nicht aus einer komplizierten Rechnung, sondern aus der Normierung der Ladungsverteilung.
    \item Eine homogen geladene Kugel hat einen oszillierenden Formfaktor, weil die Ladungsverteilung am Rand abrupt endet.
    \item Rutherford-Streuung ist Coulomb-Streuung. Sie versagt, wenn die Kerne so nahe kommen, dass Kernkr\"afte und endliche Kerngr\"o\ss{}e wichtig werden.
    \item Gr\"o\ss{}ere kinetische Energie bedeutet kleinere Ann\"aherungsdistanz. Gr\"o\ss{}erer Streuwinkel bedeutet ebenfalls kleinere Ann\"aherungsdistanz.
\end{itemize}
\end{summarybox}

\begin{formulabox}[title={Formeln, die man f\"ur dieses Blatt k\"onnen sollte}]
\begin{align*}
    F(\vect{q})&=\int e^{i\vect{q}\cdot\vect{x}/\hbar}f(\vect{x})\dd^3x,\\
    F(q)&=4\pi\int_0^\infty\frac{\sin(qr/\hbar)}{qr/\hbar}f(r)r^2\dd r,\\
    f_{\mathrm{hom}}(r)&=\begin{cases}
        \dfrac{3}{4\pi R^3}, & r\le R,\\
        0, & r>R,
    \end{cases}\\
    F_{\mathrm{hom}}(q)&=3\frac{\sin u-u\cos u}{u^3},\qquad u=\frac{qR}{\hbar},\\
    R&=r_0A^{1/3},\\
    \delta&=\frac{\delta_0}{2}\paren{1+\frac{1}{\sin(\theta/2)}},\\
    \delta_0&=\frac{Z_1Z_2\cdot\SI{1.44}{MeV.fm}}{E_{\mathrm{kin}}}.
\end{align*}
\end{formulabox}

\begin{notebox}[title=Pr\"ufungsintuition]
Wenn in einer Aufgabe nach einem Formfaktor gefragt wird, zuerst die Symmetrie der Dichte ansehen. Wenn sie kugelsymmetrisch ist, sofort Achse entlang $\vect{q}$ legen. Wenn in einer Rutherford-Aufgabe nach Kerngr\"o\ss{}eninformation gefragt wird, immer den kleinsten Abstand mit dem Kernradius vergleichen.
\end{notebox}

\end{document}
